\section{A sixteen-integer equation with a smaller arithmetic circuit}
\label{sec:signed16}

This section gives an explicit intermediate tradeoff between the direct
Pell/compiler equations and the smaller-variable construction elsewhere in
this report.  The equation has sixteen unrestricted integer unknowns,
degree
\[
  813056\cdot5^{32}+10016
  =18930435180664062500000010016,
\]
and an evaluator with \(1617\) binary arithmetic operations.  It combines
package 30's 31-variable quartic, package 29's signed Jones interface and
relation-combining construction, and package 27's digit-polynomial
evaluation identities \cite{repo27,repo29,repo30}.  The new steps are a cofinal restriction of the
coding base, the removal of its sign guard, smaller root-separation weights,
and an automatic nonzero condition.  No minimum-variable or global
operation-minimality claim is made.

Every abbreviation below denotes a polynomial expression, not another
quantified unknown.  Every exponent in the equation is a fixed natural
number.  The large fixed coefficients are defined by finite recipes; they
are not existential parameters.

\subsection{The fixed instance and the upper-bound coefficient}

Let \(Q_{31}\) be the sum of the 22 squared quadratic residuals in the
\emph{unshifted-transport} parent
\texttt{signed16/data/parent\_quartic31.json}.  This is the package-30 parent,
with its whole four-port compiler tuple instantiated by the
\(121\)-by-\(8801\) construction of this report.  In particular, this parent
is not silently replaced by the locally shifted quartic evaluator.
The coefficient interface is the same polynomial interface in
\((\beta,K_0,M_C,M_F)\); all four ports use this one complete tuple.

Put \(\nu=31\), \(L=5^{32}\), and define the normalized polynomial
\[
 {\cal P}(x,n_1,\ldots,n_\nu)
       =Q_{31}(n_1+1,\ldots,n_\nu+1).
\]
The dummy input \(x\) is unused.  The degree is four, and
\({\cal P}(x,0,\ldots,0)>0\): the fifth parent residual is \(1\) when all
positive parent coordinates equal \(1\).
Write
\begin{align*}
 {\cal P}(x,n_1,\ldots,n_\nu)
    &=\sum_{\substack{\alpha\in\N^{\nu+1}\\|\alpha|\leq4}}
       p_\alpha (x,n_1,\ldots,n_\nu)^\alpha,\\
 c_\alpha&=(4-|\alpha|)!\left(\prod_{j=0}^{\nu}\alpha_j!\right)p_\alpha,
 \qquad
 \operatorname{exp}(\alpha)=\sum_{j=0}^{\nu}\alpha_j5^j .
\end{align*}
There are 402 nonzero symbolic weighted-coefficient records.  The omitted
records mean \(c_\alpha=0\), not omitted baseline digits.  The complete
finite coefficient bounds supplied with the new tuple establish
\[
 |c_\alpha|+1<2^{4269855}\qquad (|\alpha|\leq4).
\]
Set
\begin{align}
 z_*&=2^{4269855},&
 e_*(T)&=\sum_{|\alpha|\leq4}
        (z_*+c_\alpha)T^{L-\operatorname{exp}(\alpha)},\notag\\
 \upsilon_*&=\sum_{j=1}^{\nu}(2z_*)^{5^j},&
 y_*&=e_*(2z_*),&
 \kappa&=(2z_*)^{L+1}.
 \label{eq:s16-index}
\end{align}
Thus, exactly,
\[
 \kappa
 =2^{4269856(5^{32}+1)}
 =2^{99415332078933715820316769856}.
\]
The final polynomial uses \(z_*\), \(\kappa\), and the weighted coefficients
as fixed leaves.  Neither \(y_*\) nor \(\upsilon_*\) is a variable or a leaf
needed by its evaluator.

\begin{lemma}[The fixed upper bound]\label{lem:s16-kappa}
The recipes in \eqref{eq:s16-index} define positive integers and
\(y_*<\kappa\).  They give the Jones index for the normalized parent
\({\cal P}\) with parameter \(\nu=31\).
\end{lemma}
\begin{proof}
Each digit \(z_*+c_\alpha\) lies strictly between \(1\) and \(2z_*\).
Since every \(\alpha_j\leq4\), uniqueness of base-five expansions makes
the positions \(\operatorname{exp}(\alpha)\) distinct.  Also
\(\operatorname{exp}(\alpha)\leq4\cdot5^\nu<L\).
Consequently \(y_*\) is a base-\(2z_*\) integer whose positions lie in
\(\{0,\ldots,L\}\), with each digit less than \(2z_*\).  Therefore
\[
 0<y_*
 \leq(2z_*-1)\sum_{r=0}^{L}(2z_*)^r
 =(2z_*)^{L+1}-1<\kappa.
\]
The coefficient inequality, the power-of-two choice of \(z_*\), and the
displayed definitions of \(\upsilon_*,y_*\) are exactly the admissible
index conditions used by the inherited parametric Jones theorem.  This
uses that theorem at \(\nu=31\), not the convenience existence theorem
whose internal parent has 58 variables.
\end{proof}

\subsection{The literal sixteen-coordinate equation}

The unrestricted integer coordinates, in the emitted order, are
\begin{equation}
 t,g_0,h_0,s_0,w_0,\phi,i_0,j_0,\quad
 z_1,z_2,z_3,z_4,z_5,z_6,\quad v,u .
 \label{eq:s16-coordinates}
\end{equation}
First define
\[
 b=\kappa t^2,\qquad
 g=g_0+1,\quad h=h_0+1,\quad s=s_0+1,\quad
 w=w_0+1,\quad i=i_0+1,\quad j=j_0+1,
\]
and the fixed-degree coding polynomials
\begin{gather}
 {\cal B}=b^5,\quad q={\cal B}^{L},\quad
 \ell=\sum_{r=1}^{\nu}{\cal B}^{5^r},\quad
 \lambda=\sum_{r=0}^{4L-1}{\cal B}^{r},\quad
 e=e_*({\cal B}),\quad \theta={\cal B}-2z_*,\notag\\
 {\cal N}=q^{16},\qquad G_{\rm cap}=32b{\cal B}^{5^\nu}.
 \label{eq:s16-coding}
\end{gather}
The packed rank is
\begin{align}
 R={}&
 \left[
 g+eq^3+\ell q^5+
 \left\{
  2(e-z_*\lambda)(1+{\cal B}+g)^4+
  \lambda{\cal B}+\lambda{\cal B}q^4
 \right\}q^7
 \right]({\cal N}^2-{\cal N})\notag\\
 &+
 \left[
 q^3-b\ell+\ell+\theta\lambda q^3+({\cal B}-2)q^8
 \right]({\cal N}^2-1).
 \label{eq:s16-rank}
\end{align}
Continue with the unchanged Jones ratio expressions:
\begin{align}
 U&={\cal N}^2w,& Y&={\cal N}^2s,& M&=RY,&
 P&=2M^2U,\notag\\
 K&=R+1+h(P-1),&
 A&=M(U+1),&
 B_{\rm J}&=2R+1,& W_{\rm J}&=bw,\notag\\
 C&=B_{\rm J}+W_{\rm J}+\phi,&
 J&=4A-5,&
 Q&=3W_{\rm J}C-2(W_{\rm J}^2-1),\notag\\
 D&=(A^2-1)C^2+1,&
 E&=iJDC^2,&
 F&=(A^2-1)E^2+1,\notag\\
 G_1&=A+F(CD+1-A),&
 H&=B_{\rm J}+jC,&
 I&=(G_1^2-1)H^2+1.
 \label{eq:s16-ratio-block}
\end{align}
Finally put
\begin{align}
 a_0&=(P^2-1)K^2+1,&
 a_1&=DFI,\notag\\
 m&=K^2G_{\rm cap}
       -4(C-KY)^2G_{\rm cap}-8gK^2,&
 D_0&=JF,&
 T_0&=QF+J(H-C).
 \label{eq:s16-tests}
\end{align}
The auxiliary letter \(m\) denotes the margin in this section.

The five base coordinates that retain explicit sign guards are
\[
 (x_1,x_2,x_3,x_4,x_5)=(g_0,s_0,w_0,\phi,j_0).
\]
Starting with \(C_0=1\), form the six-stage stack
\begin{align}
 d_r&=3x_r+2,&
 B_r&=1+4d_rC_{r-1}^2z_r^2,&
 C_r&=C_{r-1}B_r &&(1\leq r\leq5),\notag\\
 d_6&=3m-1,&
 B_6&=1+4d_6C_5^2z_6^2,&
 C_*&=C_6=C_5B_6 .
 \label{eq:s16-stack}
\end{align}
The circuit names the stack outputs \texttt{guard1}, \ldots,
\texttt{guard6}; the ratio-block expression \(B_{\rm J}=2R+1\)
is separate.

For arbitrary integer arguments define
\begin{equation}
 {\cal R}(D',T',v)=
 \big((T'+D'v)^3-D'^2(T'+D'v)-D'^3\big)
 \big((T'+D'v)^2+23D'^2\big).
 \label{eq:s16-homogeneous}
\end{equation}
Using this polynomial, set
\begin{align}
 Z&=i\prod_{r=1}^{6}z_r,&
 S&=ZD_0^5,& T&={\cal R}(D_0,T_0,v),\notag\\
 w_1&=C_*,&
 V_0&=P^2+K^2+2,&
 w_2&=C_*V_0,\notag\\
 X&=Su-T,&
 A_0&=S^2C_*,&
 A_1&=S^2w_1^2a_0,&
 A_2&=S^2w_2^2a_1,\notag\\
 V&=X^2+A_2-A_0-A_1,&
 W&=V^2+4X^2A_2-4A_0A_1 .
 \label{eq:s16-finalizer}
\end{align}
The promised ordinary integer polynomial is
\begin{equation}
 \boxed{\displaystyle F_{16}=W^2-16V^2X^2A_2.}
 \label{eq:s16-polynomial}
\end{equation}
In particular, \(t\) is deliberately absent from \(Z\).  Its nonvanishing
will follow from the recovered margin; omitting it saves one multiplication.

\subsection{The exact imported interface and the new cofinal restriction}

For an integer \(a\), write
\(\operatorname{Sq}(a)\) for the assertion
\(\exists r\in\Z\;(a=r^2)\).
The inherited Jones ratio equivalence applies to the eight
\emph{nonnegative} coordinates
\[
 \varepsilon,g_0,h_0,s_0,w_0,\phi,i_0,j_0,
 \qquad b=y_*+1+\varepsilon.
\]
For the fixed normalized parent, it says that
\begin{equation}
 \operatorname{Sq}(a_0),\quad \operatorname{Sq}(a_1),\quad
 D_0\mid T_0,\quad m>0
 \label{eq:s16-jr}
\end{equation}
are jointly solvable if and only if \({\cal P}\) has a nonnegative
31-tuple zero.  The ratio expressions in
\eqref{eq:s16-coding}--\eqref{eq:s16-tests} are precisely those in this
interface.  This section does not apply a signed interpretation to an
unspecified nine-variable polynomial.

The following strengthened completeness fact is needed to replace
\(b=y_*+1+\varepsilon\) by \(b=\kappa t^2\).

\begin{lemma}[Cofinal powers of two]\label{lem:s16-cofinal}
Given a nonnegative parent witness
\({\cal P}(1,n_1,\ldots,n_\nu)=0\), the Jones ratio conditions
\eqref{eq:s16-jr} have witnesses with \(b=2^a\) for every sufficiently
large integer \(a\).  In particular, they have witnesses with
\(b=\kappa t^2\) and \(t>0\).
\end{lemma}
\begin{proof}
Let \(\xi=(1,n_1,\ldots,n_\nu)\).
The constructive master-system proof
\texttt{Master3.USys\_of\_mem} chooses
\[
 a_0'=y_*+\sum_{r=0}^{\nu}\xi_r+1,\qquad b=2^{a_0'}.
\]
Inspection of that complete proof shows that the equality specifying
\(a_0'\) is used only to obtain
\[
 a\geq1,\qquad y_*<b,\qquad \xi_r<b\quad(0\leq r\leq\nu).
 \tag{C}
\]
These conclusions hold for every \(a\geq a_0'\), because
\(a<2^a\).  The rest of the construction depends on (C), on \(b=2^a\),
and on the given polynomial zero; it imposes no upper bound on \(a\).

More explicitly, at this chosen base one takes the polynomial evaluations
\[
 e=e_*({\cal B}),\qquad
 \ell=\sum_{r=1}^{\nu}{\cal B}^{5^r},\qquad
 g=\sum_{r=1}^{\nu}\xi_r{\cal B}^{5^r},\qquad {\cal B}=b^5,
\]
together with \(q,\lambda,\theta\) from \eqref{eq:s16-coding}.
Normalization of the parent implies \(g>0\).
The digit bounds and the source's carry-free coding identity give
\[
 e<2z_*{\cal B}^{L},\qquad
 \ell<2{\cal B}^{5^\nu},\qquad
 g<b{\cal B}^{5^\nu}.
\]
The necessary positive slack is
\[
 \alpha_{\rm master}
       =(b-y_*)q^2-e\ell g^2.
\]
It is positive uniformly in the larger base: putting \(r=5^\nu\),
\[
 e\ell g^2
 <4z_*b^2{\cal B}^{L+3r}
 \leq{\cal B}^{L+3r+1}
 \leq{\cal B}^{2L}=q^2.
\]
Here \(4z_*b^2\leq b^5={\cal B}\) follows from \(b>y_*\geq2z_*\),
and \(L+3r+1\leq2L\) follows from \(L=5r\).
The two master quotients
\[
 \frac{e-y_*}{\theta},\qquad
 \frac{\ell-\upsilon_*}{\theta}
\]
are positive integers: divisibility is polynomial evaluation modulo
\({\cal B}-2z_*\), and positivity follows from
\({\cal B}>2z_*\).  The three bit-disjointness identities in the master
system are the inherited Jones coding identities; their hypotheses are
(C), \(b=2^a\), and the unchanged parent zero.  Thus the same explicit
construction works at every such larger base.

The proof of \texttt{NineReduce.reduced\_of\_mem} keeps this very base
and gives the required packed-rank/binomial condition, as well as
\(g<b{\cal B}^{5^\nu}\).  The theorem
\texttt{NineRatio.ratio\_complete} then takes \(b\) as a fixed parameter;
its hypothesis \(32g<G_{\rm cap}\) follows from that last inequality.
It supplies the positive \(h,s,w,i,j\) and nonnegative \(\phi\) in
\eqref{eq:s16-jr}, without selecting a different base.

Write \(\kappa=2^c\), where
\(c=4269856(L+1)\).  Choose any \(r\geq0\) with
\(c+2r\geq a_0'\) and set \(t=2^r\).  Then
\[
 \kappa t^2=2^{c+2r}
\]
is one of the permitted bases.  This gives the required positive \(t\)
and proves the assertion constructively.
\end{proof}

\subsection{Signed coordinates whose guards are unnecessary}

\begin{lemma}[A denominator that never vanishes]\label{lem:s16-denominator}
For every integer assignment of the eight base coordinates,
\(D_0=JF\neq0\).
\end{lemma}
\begin{proof}
The integer \(J=4A-5\) is nonzero.  If
\(F=(A^2-1)E^2+1=0\), then
\((1-A^2)E^2=1\), so \(A=0\) and \(E=\pm1\).
But \(J=-5\) when \(A=0\), and \(E=iJDC^2\) is divisible by \(5\).
This is impossible.
\end{proof}

\begin{lemma}[The sign of \(i\)]\label{lem:s16-isign}
If \(i\neq0\), replacing \(i\) by \(|i|\) preserves
\(a_0,a_1,D_0,T_0,m\), and makes \(i_0=|i|-1\) nonnegative.
\end{lemma}
\begin{proof}
The only unsquared use of \(i\) is in \(E=iJDC^2\).
Among the five expressions under consideration, every dependence on
\(E\) passes through \(E^2\) in \(F\).  Thus they depend only on \(i^2\).
\end{proof}

\begin{lemma}[The \(h\) guard follows from the margin]\label{lem:s16-hguard}
Suppose \(b>y_*\), the five coordinates
\(g_0,s_0,w_0,\phi,j_0\) are nonnegative, \(a_0\) is an integer square,
and \(m>0\).  Then \(h\geq1\), hence \(h_0\geq0\).
\end{lemma}
\begin{proof}
We have \(g,s,w,j\geq1\), \(\phi\geq0\),
\(b\geq2\), \({\cal B}=b^5>2z_*\), \(e\geq0\),
and \(G_{\rm cap}>0\).
The margin inequality first gives \(K\neq0\) and
\begin{equation}
 g<G_{\rm cap}/8=4b{\cal B}^{5^\nu},
 \qquad
 4(C-KY)^2<K^2.
 \label{eq:s16-margin-consequences}
\end{equation}
Indeed \(m\leq K^2(G_{\rm cap}-8g)\); then discarding the strictly
positive term \(8gK^2\) gives the second inequality.

Put \(r=5^\nu\).  Since \({\cal B}>5\) and \(b\leq{\cal B}\),
\[
 1+{\cal B}+g
 <1+{\cal B}+4{\cal B}^{r+1}
 \leq5{\cal B}^{r+1}
 <{\cal B}^{r+2}
 \leq{\cal B}^{5r}=q.
\]
Also \(\ell<{\cal B}^{r+1}\), hence
\(b\ell<{\cal B}^{r+2}\leq q<q^3\).
The potentially negative part of the first bracket in
\eqref{eq:s16-rank} is absorbed by its positive term:
\[
 \lambda{\cal B}q^4
       >2z_*\lambda(1+{\cal B}+g)^4 .
\]
All the remaining terms of that bracket are nonnegative, and \(g>0\).
The second bracket is positive because \(q^3>b\ell\), \(\theta>0\),
and \({\cal B}>2\).
It follows that
\[
 R\geq{\cal N}^2-{\cal N}\geq2.
\]
Thus \(U,Y\geq1\), \(C=2R+1+bw+\phi>0\), and
\[
 P=2R^2Y^2U\geq2R^2>R+2.
\]

The second margin consequence rules out \(K\leq0\):
\(K=0\) is impossible, and \(K<0\) would give
\(|C-KY|=C+|K|Y>|K|\).
Therefore \(K>0\).
If \(h\leq0\), the relation \(K=R+1+h(P-1)\) gives
\(1\leq K\leq R+1<P\).
For \(K=1\) it would imply \(P-1\mid R\), contradicting
\(0<R<P-1\).
For \(2\leq K<P\),
\[
 (PK-1)^2<(P^2-1)K^2+1<(PK)^2,
\]
since the two differences are \(K(2P-K)>0\) and \(K^2-1>0\).
This contradicts squarehood of \(a_0\).
Consequently \(h\geq1\).
\end{proof}

\subsection{The coprime Pell stack}

\begin{lemma}[Six guards in one square test]\label{lem:s16-stack}
Subject to \(z_1\cdots z_6\neq0\), squarehood of \(C_*\) in
\eqref{eq:s16-stack} implies
\begin{equation}
 g_0,s_0,w_0,\phi,j_0\geq0,\qquad m>0.
 \label{eq:s16-guards}
\end{equation}
Conversely, if these inequalities hold, the six \(z_r\) can be chosen
nonzero so that \(C_*\) is a square.
\end{lemma}
\begin{proof}
Every \(B_r\equiv1\pmod4\), so no \(B_r\) vanishes.
Also \(\gcd(B_r,C_{r-1})=1\).
Inductively, the six factors \(B_r\) are pairwise coprime in absolute
value.  If their product is a square, unique prime factorization shows
that every \(|B_r|\) is a square.  A negative square is \(0\) or \(3\)
modulo \(4\), so \(B_r\equiv1\pmod4\) forces every factor to be a
positive square.  If any \(d_r<0\), the nonzero integers
\(C_{r-1},z_r\) would give
\(B_r=1+4d_rC_{r-1}^2z_r^2\leq-3\), a contradiction.
Finally \(d_r\equiv2\pmod3\), so it cannot be zero.
Thus \(d_r>0\) for all \(r\).  For \(r\leq5\) this means \(x_r\geq0\);
for \(r=6\) it means \(m\geq1\).

For the converse, each \(d_r>0\) and \(d_r\equiv2\pmod3\), so it is
not a square.  The Pell equation
\[
 a^2-d_r v^2=1
\]
has a nontrivial positive solution.  It has a positive solution with
\(v\) divisible by any prescribed positive integer \(M\): the matrix
\(\left(\begin{smallmatrix}a&d_rv\\v&a\end{smallmatrix}\right)\)
has determinant \(1\), and some positive power of its reduction modulo
\(M\) is the identity in the finite group
\(\mathrm{SL}_2(\Z/M\Z)\).
Choose \(M=2|C_{r-1}|\), and put \(z_r=v/(2C_{r-1})\).
Then \(z_r\) is a nonzero integer and \(B_r=a^2\).
Starting with \(C_0=1\), induction makes every \(C_r\) a positive square.
\end{proof}

A useful strengthening is available before any guard inequality is known.
If all \(z_r\neq0\), then \(d_r\neq0\), \(C_{r-1}\neq0\), and hence
\(B_r\neq1\).  The congruence \(B_r\equiv1\pmod4\) also excludes
\(B_r=-1\).  Therefore
\begin{equation}
 |C_*|>1
 \quad\hbox{whenever}\quad z_1\cdots z_6\neq0.
 \label{eq:s16-stack-large}
\end{equation}
This preguard estimate is what permits the smaller weights.

\subsection{Divisibility and nonvanishing in one relation}

\begin{lemma}[An intersective polynomial]\label{lem:s16-intersective}
The polynomial
\[
 \Pi(T)=(T^3-T-1)(T^2+23)
\]
has no rational zero, but has an integer zero modulo every positive
integer modulus.
\end{lemma}
\begin{proof}
The only possible rational roots of the monic cubic are \(1,-1\),
and neither is a root.  The quadratic has no real root.

For an odd prime \(p\neq23\), the cubic has discriminant \(-23\).
If it has a root modulo \(p\), the root is simple and lifts to every
power of \(p\).
If it has no such root, it is irreducible of degree three.
Frobenius cycles its three roots in \(\mathbb F_{p^3}\), an even
permutation, so their Vandermonde product lies in \(\mathbb F_p\).
Its square is \(-23\), which is nonzero.  Thus \(T^2+23\) has a
simple root modulo \(p\), again lifting to every power.
At \(p=23\), the cubic has the simple root \(3\), since
\(3^3-3-1=23\) and \(3(3^2)-1=26\not\equiv0\pmod{23}\).

For \(p=2\), use \(-23\equiv1\pmod8\).
If an odd \(a\) is a root of \(T^2+23\) modulo \(2^e\), \(e\geq3\),
one of \(a\) and \(a+2^{e-1}\) is a root modulo \(2^{e+1}\):
their squares differ by \(2^ea\) modulo \(2^{e+1}\).
This supplies roots for all powers of two.
The Chinese remainder theorem combines roots for the prime-power
factors of any positive modulus.
\end{proof}

\begin{lemma}[Homogeneous divisibility merger]\label{lem:s16-merger}
For integers \(D',T',Z'\) with \(D'\neq0\),
\begin{equation}
 \big(D'\mid T'\ \hbox{and}\ Z'\neq0\big)
 \quad\Longleftrightarrow\quad
 \exists v\in\Z:\ Z'D'^5\mid{\cal R}(D',T',v).
 \label{eq:s16-merger}
\end{equation}
Moreover \({\cal R}(D',T',v)\neq0\) for every integer \(v\).
\end{lemma}
\begin{proof}
The rational identity
\[
 {\cal R}(D',T',v)=D'^5\Pi(T'/D'+v)
\]
proves nonvanishing by Lemma~\ref{lem:s16-intersective}.
If \(T'=D'q\) and \(Z'\neq0\), choose \(v\) so that
\(Z'\mid\Pi(q+v)\); this proves the forward implication.

Conversely, \(Z'=0\) would make the displayed divisibility assert
\({\cal R}=0\), which is impossible.
If \(D'\nmid T'\), there is a prime \(p\) with
\(\alpha=v_p(|D'|)>\beta=v_p(|T'|)\).
Here \(T'\neq0\).
For \(r=T'+D'v\), one has \(v_p(r)=\beta\).
In \(r^3-D'^2r-D'^3\), the first term has the unique least valuation
\(3\beta\); in \(r^2+23D'^2\), it has the unique least valuation
\(2\beta\).  Hence
\[
 v_p({\cal R})=5\beta<5\alpha=v_p(D'^5),
\]
contradicting the asserted divisibility.  Signs of \(D'\) and \(Z'\)
do not affect this argument.
\end{proof}

\subsection{The three-root norm with signed weights}

\begin{lemma}[The finalizer]\label{lem:s16-finalizer}
Fix all coordinates except \(u\).
For the expressions above, \(\exists u\in\Z:F_{16}=0\) is equivalent to
\begin{equation}
 S\neq0,\qquad S\mid T,\qquad
 \operatorname{Sq}(C_*),\quad
 \operatorname{Sq}(a_0),\quad
 \operatorname{Sq}(a_1).
 \label{eq:s16-finalizer-interface}
\end{equation}
\end{lemma}
\begin{proof}
Lemma~\ref{lem:s16-denominator} gives \(D_0\neq0\), and
Lemma~\ref{lem:s16-merger} gives \(T\neq0\).
If \(S=0\), then \(A_0=A_1=A_2=0\), \(X=-T\), and
\(F_{16}=T^8\neq0\).
Thus every zero has \(S\neq0\).
Since \(S=i(z_1\cdots z_6)D_0^5\), every \(z_r\) and \(i\) is
nonzero; in particular \eqref{eq:s16-stack-large} applies.

For all real \(P,K\),
\begin{equation}
 \sqrt{|(P^2-1)K^2+1|}\leq P^2+K^2+1 .
 \label{eq:s16-a0-bound}
\end{equation}
Indeed
\(|(P^2-1)K^2+1|\leq(P^2+1)K^2+1\), and the difference between
\((P^2+K^2+1)^2\) and this last upper bound is
\[
 P^4+K^4+P^2K^2+2P^2+K^2\geq0.
\]
It follows from \(|C_*|>1\) that the possibly signed weights satisfy
\begin{align}
 |w_1|&>\sqrt{|C_*|},\notag\\
 |w_2|&>\sqrt{|C_*|}+|w_1|\sqrt{|a_0|}.
 \label{eq:s16-separation}
\end{align}
These inequalities require neither positive radicands nor recovered
guard inequalities.

Choose complex roots
\(\rho_0^2=C_*,\rho_1^2=a_0,\rho_2^2=a_1\).
Algebraic multiplication gives
\begin{equation}
 F_{16}
  =\prod_{\epsilon_0,\epsilon_1,\epsilon_2\in\{-1,1\}}
     \left[
       X-S(\epsilon_0\rho_0+
             \epsilon_1w_1\rho_1+\epsilon_2w_2\rho_2)
     \right].
 \label{eq:s16-eight-factor}
\end{equation}
For example, first multiplying the \(\rho_0,\rho_1\) signs gives
\(\big((X\mp\sqrt{A_2})^2-A_0-A_1\big)^2-4A_0A_1\);
the two remaining factors are
\(W\mp4VX\sqrt{A_2}\), whose product is
\eqref{eq:s16-polynomial}.

A zero makes one of the signed sums in
\eqref{eq:s16-eight-factor} equal to \(X/S=u-T/S\), a rational number.
Distinct sign patterns on nonzero roots give distinct sums:
at the largest index where signs differ, the absolute value of the
weighted root is at least the absolute value of its weight, since a
nonzero integer radicand has root magnitude at least \(1\).
Inequalities \eqref{eq:s16-separation} make this strictly larger than
the sum of the magnitudes of all preceding weighted roots.
The triangle inequality therefore precludes cancellation.  This remains
valid for imaginary roots of negative radicands and for negative weights.
Zero roots may simply be omitted from this comparison.

The field \(\mathbb Q(\rho_0,\rho_1,\rho_2)\) is the splitting field of
three quadratic polynomials, hence a finite Galois extension of
\(\mathbb Q\).  Each automorphism changes root signs.
Since the encoded sum is rational, every automorphism fixes it.
The separation just proved forces every nonzero root in it to be fixed
individually.  Therefore all three roots are rational.
A rational root of a monic polynomial \(T^2-a\), \(a\in\Z\), is an
integer.  Thus \(C_*,a_0,a_1\) are nonnegative integer squares.
The encoded sum is now an integer, so \(T/S\) is an integer as well.
This proves necessity of \eqref{eq:s16-finalizer-interface}.

Conversely, under that interface choose integer square roots and set
\[
 u=T/S+\rho_0+w_1\rho_1+w_2\rho_2 .
\]
This is an integer and makes a factor in
\eqref{eq:s16-eight-factor} vanish.
\end{proof}

\subsection{Soundness, completeness, and the consistency equivalence}

\begin{theorem}[The sixteen-integer presentation]\label{thm:s16}
For the fixed coefficients above, ordinary mathematics formalizable
in ZFC proves
\[
 \big(\exists {\bf x}\in\Z^{16}:F_{16}({\bf x})=0\big)
 \ \longleftrightarrow\
 \big(\text{there is a finite syntactic ZFC refutation}\big).
\]
Consequently,
\[
 \mathrm{ZFC}\ \vdash\
 \big(\nexists{\bf x}\in\Z^{16}:F_{16}({\bf x})=0\big)
 \ \longleftrightarrow\ \operatorname{Con}(\mathrm{ZFC}).
\]
\end{theorem}
\begin{proof}
For soundness, let \(F_{16}=0\).
Lemma~\ref{lem:s16-finalizer} gives \(S\neq0\), \(S\mid T\), and
the three square conditions.
Lemma~\ref{lem:s16-merger} gives \(D_0\mid T_0\) and
\(i z_1\cdots z_6\neq0\).
Lemma~\ref{lem:s16-stack} recovers the five nonnegativity guards
and \(m>0\).

At this point \(t\neq0\) follows without a separate nonzero factor.
If \(t=0\), then \(b=\kappa t^2=0\) and \(G_{\rm cap}=0\), while
\(g=g_0+1\geq1\).  The literal margin formula would give
\[
 m=-8gK^2\leq0,
\]
a contradiction.
Hence \(b=\kappa t^2\geq\kappa>y_*\), and
\(\varepsilon=b-y_*-1\) is a nonnegative integer.
Lemma~\ref{lem:s16-hguard} gives \(h_0\geq0\).
Lemma~\ref{lem:s16-isign} replaces \(i\) by \(|i|\), preserving all
four Jones conditions and making \(i_0\geq0\).
We have therefore recovered an admissible nonnegative tuple for the
unchanged Jones interface \eqref{eq:s16-jr}.
The inherited ratio soundness theorem gives a zero of the parent
quartic.  Its complete native/Pell compiler interface gives a blank
K215 halting computation, and the calibrated K215 theorem gives a
finite syntactic ZFC refutation.

For completeness, a finite ZFC refutation gives a blank K215 halting
run and then a positive parent-quartic solution by those same two
interfaces.
After the coordinate shift to the normalized parent, apply
Lemma~\ref{lem:s16-cofinal}.
It gives a Jones tuple satisfying \eqref{eq:s16-jr} at
\(b=\kappa t^2\) with \(t>0\), all the ordinary base coordinates
nonnegative, and \(i>0\).
Lemma~\ref{lem:s16-stack} supplies nonzero \(z_1,\ldots,z_6\)
for which \(C_*\) is a square.
Now \(Z=i z_1\cdots z_6\neq0\); the inherited Jones tuple has
\(D_0\mid T_0\).
Lemma~\ref{lem:s16-merger} chooses \(v\) with \(S\mid T\).
The three radicands are squares, so
Lemma~\ref{lem:s16-finalizer} chooses the final integer \(u\)
and gives \(F_{16}=0\).

Every extra root, quotient, or exponent selected in this proof is a
proof witness used to establish equivalence.  The displayed equation
itself quantifies only the sixteen coordinates in
\eqref{eq:s16-coordinates}.
Negating the existence equivalence gives the stated consistency
equivalence.  Neither direction assumes or proves
\(\operatorname{Con}(\mathrm{ZFC})\).
\end{proof}

\subsection{Degree and arithmetic size}

All degrees in this subsection give each of the sixteen unknowns
degree one and every fixed coefficient degree zero.
The largest term in the packed rank is
\(\lambda{\cal B}q^{11}{\cal N}^2\);
it gives \(\deg R=470L\).
The largest term in the margin is
\(-4K^2Y^2G_{\rm cap}\), giving \(\deg m=4442L+12\).
The useful intermediate degrees are
\[
\begin{array}{c|r@{\qquad}c|r}
 R&470L&P&1900L+3\\
 K&1900L+4&A&1110L+2\\
 a_0&7600L+14&a_1&49280L+68\\
 m&4442L+12&C_5&363\\
 C_*&4442L+1103&D_0&13750L+20\\
 S&68750L+107&T&68750L+105\\
 X&68750L+108&w_2&8242L+1111\\
 A_0&141942L+1317&A_1&153984L+2434\\
 A_2&203264L+2504&&
\end{array}
\]
For the first five stack stages the degree recurrence is
\(c_r=3c_{r-1}+3\), starting at \(c_0=0\).
The last stage has degree
\(3\cdot363+(4442L+12)+2=4442L+1103\).
The seven factors of \(Z=i\prod_{r=1}^6z_r\) account for the constant
\(107=5\cdot20+7\) in the degree of \(S\).

The term \(A_2\) has strictly greater degree than \(X^2,A_0,A_1\).
Thus \(V\) has the leading form of \(A_2\), \(W\) that of \(A_2^2\),
and the final polynomial that of \(A_2^4\).
All of these leading forms are nonzero: the rank has the displayed
unique top term, the margin has its unique negative top term, and
subsequent relevant leading forms are products and squares of
nonzero forms.  In particular, after the authentic positive fixed
coefficients are substituted,
\[
 \deg F_{16}=4(203264L+2504)=813056L+10016.
\]
This is an exact degree assertion, not just a structural upper bound
that ignores cancellations.

The literal pruned evaluator contains
\[
 697\ \hbox{additions}
 \;+\;26\ \hbox{subtractions}
 \;+\;894\ \hbox{multiplications}
 \;=\;1617\ \hbox{operations}.
\]
Each binary operation, including a square and multiplication by a
fixed numeral, is charged.  Reuse and fixed leaves are free, consistently
with the repository's arithmetic-circuit convention.  Fixed exponentiation
is expanded into multiplication gates; no exponentiation oracle is used.

For clarity, the compact evaluation of the enormous digit baseline
does not omit its zero-coefficient positions.  If
\(p_r={\cal B}^{5^r}\), then the baseline is represented through the
complete homogeneous polynomial of degree four, and
\[
 \lambda
 =\left(\prod_{r=0}^{\nu}
       (1+p_r+p_r^2+p_r^3+p_r^4)\right)(1+q)(1+q^2).
\]
These are polynomial identities.
The full baseline contains
\(\binom{36}{4}=58905\) distinct positions; the 402 nonzero
weighted-coefficient records supply its sparse correction.

\subsection{Verification boundary}

The final binding is
\texttt{signed16/data/INSTANCE.json}; it names the whole compiler
tuple, parent, weighted-coefficient records, and index recipes.
The literal division-free polynomial is
\texttt{signed16/generated/signed16.json}.
The new reduction does not claim a fresh end-to-end Lean or Rocq
certificate.  The source Lean files explain the precise Jones interfaces
and the cofinal construction, while the native compiler and its K215
calibration remain identified mathematical proof dependencies.

The independent data-only auditor imports neither the emitter nor its DAG
utility.  It checks all four decompressed port hashes, enumerates the
58905 baseline positions directly, evaluates the geometric sum by a
separate method, and compares 288 named-node values on six signed
assignments over two primes, including \(t=0\).
It also checks the eight-factor norm in the rank-eight ring with three
formal square roots on 100 exact assignments.
For the final numerical instance, the top-degree diagonal coefficients
are nonzero modulo two primes:
\[
 121738\pmod{1000003},\qquad
 73024\pmod{1000033}.
\]
Together with the propagated degree upper bound, either nonzero residue
is an independent certificate that the top degree survives on this fixed
coefficient slice.

These checks establish transcription, coefficient binding, and algebraic
metrics.  The universal soundness and completeness claims rest on the
proofs above and the stated inherited interfaces, not on finite tests or
a search for a Diophantine solution.

